% HSK MCM/ICM LaTeX template v6.2.2
% Compile with: python scripts/render_paper.py final_latex --profile mcm_icm --clean
\documentclass[11pt,letterpaper]{article}

\usepackage[margin=1in]{geometry}
\usepackage{amsmath,amssymb,amsthm,bm}
\usepackage{graphicx}
\usepackage{float}
\usepackage{booktabs}
\usepackage{longtable}
\usepackage{multirow}
\usepackage{listings}
\usepackage{xcolor}
\usepackage{hyperref}
\usepackage{url}
\usepackage{cite}
\usepackage{caption}
\usepackage{subcaption}
\usepackage{enumitem}
\usepackage{titlesec}
\usepackage{fancyhdr}
\usepackage{lastpage}
\usepackage{newtxtext,newtxmath}

\newcommand{\ContestName}{MCM/ICM}
\newcommand{\ContestYear}{YYYY}
\newcommand{\ProblemCode}{A}
\newcommand{\TeamNumber}{XXXXXXX}
\newcommand{\PaperTitle}{TITLE OF THE PAPER}

\hypersetup{
    colorlinks=false,
    pdfborder={0 0 0}
}

\lstset{
    basicstyle=\small\ttfamily,
    breaklines=true,
    frame=single,
    numbers=left,
    numberstyle=\tiny,
    showstringspaces=false
}

\pagestyle{fancy}
\fancyhf{}
\fancyhead[L]{Team \#\TeamNumber}
\fancyhead[R]{Page \thepage\ of \pageref{LastPage}}
\renewcommand{\headrulewidth}{0.4pt}
\titlespacing*{\section}{0pt}{12pt}{6pt}
\titlespacing*{\subsection}{0pt}{10pt}{4pt}

\begin{document}

\begin{titlepage}
\begin{center}
\vspace*{1.2cm}
{\LARGE\bfseries \ContestYear\ \ContestName\ Summary Sheet}\\[0.8cm]
{\large Problem Chosen: \textbf{\ProblemCode}}\\[0.4cm]
{\large Team Control Number: \textbf{\TeamNumber}}\\[1.5cm]
{\Huge\bfseries \PaperTitle}\\[1.5cm]
\end{center}

\section*{Summary}
% Keep the summary within the current competition limit. Cover the problem,
% model architecture, quantitative results, validation, and the main decision.
Write the one-page summary here. Replace every placeholder before submission.

\vfill
\noindent\textbf{Keywords:} keyword 1; keyword 2; keyword 3; keyword 4
\end{titlepage}

\tableofcontents
\newpage

\section{Introduction}
\subsection{Problem Context}
\subsection{Problem Restatement}
\subsection{Modeling Challenges}
\subsection{Our Contributions}
% Do not translate a Chinese competition-paper outline mechanically. Explain the
% specific objects, mechanisms, constraints, and evidence used in this problem.

\section{Assumptions and Justifications}
% Retain only assumptions that materially affect the model. State the reason,
% likely failure bias, and validation method for each assumption.

\section{Notation}
\begin{table}[H]
\centering
\caption{Main notation}
\label{tab:notation}
\begin{tabular}{llll}
\toprule
Symbol & Definition & Unit & Type \\
\midrule
$x_i$ & Example decision variable & -- & Decision \\
$s_t$ & Example state variable & -- & State \\
$\theta$ & Example parameter & -- & Parameter \\
\bottomrule
\end{tabular}
\end{table}

\section{Data and Preprocessing}
\subsection{Data Scope and Alignment}
\subsection{Missing Values, Outliers, and Transformations}

\section{Model Development}
\subsection{Subproblem 1}
\subsubsection{Mechanism and Variables}
\subsubsection{Mathematical Formulation}
\subsubsection{Solution Method}
\subsubsection{Results and Interpretation}

% Duplicate the structure above only for the actual number of subproblems.

\section{Validation and Uncertainty}
% Select validation appropriate to the task: out-of-sample evaluation,
% constraint residuals, boundary checks, multi-start or multi-algorithm tests,
% confidence intervals, calibration, or spatial robustness.

\section{Sensitivity and Robustness Analysis}
% Use one-at-a-time, joint perturbation, scenario stress testing, Bootstrap,
% Sobol, or alternative-model analysis only when supported by the problem.

\section{Strengths, Limitations, and Recommendations}
\subsection{Strengths}
\subsection{Limitations}
\subsection{Actionable Recommendations}

\section{Conclusion}
% Answer each requested question directly with quantitative results and scope.

\bibliographystyle{ieeetr}
\bibliography{references}

% Enable a letter only when the current problem explicitly requires one.
% \newpage
% \section*{Letter to the Decision Maker}
% \addcontentsline{toc}{section}{Letter to the Decision Maker}

\appendix
\section{Reproduction Notes}
\section{Supplementary Tables and Figures}
\section{Code Description}
% Full code may be submitted separately. State the input, function, output, and
% reproduction order rather than filling the main paper with long code blocks.

\end{document}
